\chapter{Introduction}

Combinatorial optimization problems lie at the foundation of numerous real-world applications and theoretical investigations. At their core, these problems involve making discrete choices that collectively maximize or minimize a certain objective under a set of constraints. The scope of such problems is extraordinarily wide: from assigning frequencies in wireless networks and scheduling manufacturing operations, to designing efficient routes for logistics or solving large-scale allocation problems in cloud computing. In artificial intelligence and computer vision, tasks such as image segmentation, stereo matching, and graphical model inference can also be phrased in terms of combinatorial optimization. Despite their importance, these problems are often difficult to solve in practice. Many fall into the class of NP-hard problems, for which no known polynomial-time algorithms exist \citep{lewis_michael_1983}. As a result, researchers and practitioners frequently resort to approximate methods, heuristics, or problem-specific relaxations in order to obtain useful solutions.

The limitations of traditional approaches to combinatorial optimization are well documented. Exact solvers based on integer programming, branch-and-bound, or constraint satisfaction can handle small to medium-sized problem instances but quickly become intractable as problem size grows \citep{wolsey_integer_2020}. Approximation algorithms and heuristics, while more scalable, are often tailored to specific problem families and lack the ability to generalize to new or unseen problem structures. Furthermore, many heuristic approaches come with no guarantees on the quality of solutions, leaving practitioners to rely on empirical validation rather than theoretical assurance. This tension between the universality of combinatorial optimization problems and the specialized, limited applicability of traditional methods motivates the search for more flexible and powerful solution paradigms.

Probabilistic graphical models, and in particular \textit{Markov Random Fields (MRFs)}, offer a promising representational framework for addressing this challenge. MRFs capture the joint distribution of a set of variables by factorizing it according to the structure of an undirected graph \citep{kindermann_relation_1980}. By assigning potentials to nodes and edges, one can encode local interactions and global constraints in a manner that naturally aligns with the structure of many combinatorial problems. Importantly, solving a combinatorial optimization problem can often be recast as an inference task in an MRF: for example, the Maximum Cut problem is equivalent to finding a ground state in an Ising model \citep{barahona_computational_1982}, while satisfiability and covering problems can be translated into maximum a posteriori (MAP) inference in corresponding factor graphs.

This connection becomes particularly compelling when viewed through the lens of statistical physics. Many combinatorial problems can be modeled as the search for low-energy states of a system governed by a Boltzmann-type distribution. In such a distribution, the probability of a configuration is proportional to $\exp(-E(x)/T)$, where $E(x)$ is an energy function encoding the objective and constraints of the optimization problem, and $T$ is a temperature parameter \citep{geman_stochastic_1984,mezard_information_2009}. Markov Random Fields provide a convenient and flexible way of modeling these Boltzmann distributions, where potentials correspond to energy contributions of local structures. Thus, combinatorial solutions can be interpreted as outcomes of a random process driven by such distributions, and MRF inference becomes the natural vehicle for exploring and learning these solutions.

While graphical models provide structure, \textit{machine learning} offers adaptability and scalability. The past decade has seen remarkable progress in learning-based methods that surpass traditional hand-designed heuristics in both accuracy and efficiency across diverse domains. Deep learning, reinforcement learning, and graph neural networks (GNNs) have all contributed to a new paradigm often described as \textit{learning to optimize} \citep{bengio_machine_2018}. In this paradigm, instead of explicitly crafting algorithms for each new combinatorial problem, one trains models to implicitly capture solution strategies from data. Such models can generalize across instances of a problem, sometimes even across related problem classes, enabling rapid inference once trained. For example, reinforcement learning has been applied to discover new heuristics for the traveling salesman problem \citep{bello_neural_2016}, while GNNs have demonstrated the ability to approximate solutions to graph partitioning and covering problems by exploiting structural regularities \citep{velickovic_graph_2017, brody_how_2021, dai_learning_2017}.

The convergence of these two perspectives—graphical modeling and machine learning—motivates the central focus of this thesis. On one hand, MRFs provide a principled representation for combinatorial optimization, unifying a wide range of problems under a common probabilistic framework. On the other hand, machine learning techniques, particularly those inspired by message passing and neural architectures, offer the potential to overcome the limitations of classical inference methods. The central question driving this research can thus be framed as: \textbf{How can machine learning enhance MRF-based inference to more effectively solve combinatorial optimization problems, both in terms of scalability and solution quality?}

This thesis specifically explores this question through the study of \textit{message passing algorithms}, a family of inference methods well-suited to MRFs. Message passing is attractive because it exploits local graph structure and can be distributed across nodes and edges, making it scalable in principle. However, classical message passing methods such as loopy belief propagation often fail to converge or provide poor approximations on loopy graphs, which are common in combinatorial optimization \citep{pearl_reverend_1982,yedidia_understanding_2003}. By integrating machine learning, it becomes possible to adapt and improve the dynamics of message passing—for example, by learning message scheduling policies, adjusting update rules, or parameterizing messages through neural networks. The aim is to retain the interpretability and structure of MRF inference while incorporating the flexibility and empirical success of modern learning techniques.

The specific contribution of this thesis lies in applying these ideas to two canonical graph-theoretic optimization problems: the \textit{Maximum Independent Set (MIS)} and the \textit{Set Cover (SC)}. These problems are both fundamental in theoretical computer science and directly relevant to applications in network analysis, resource allocation, and bioinformatics. By modeling MIS and SC within MRF frameworks and applying learned message passing inference, the thesis demonstrates how machine learning can extend classical inference into a powerful solver for combinatorial problems.

The objectives of this thesis can be described as follows. First, it seeks to provide a systematic overview of how combinatorial optimization problems can be formulated as MRFs, clarifying the connections between classical optimization problems, Boltzmann-type distributions, and probabilistic inference. Second, it aims to investigate the integration of machine learning into message passing algorithms, focusing on how learned components can improve both convergence behavior and solution quality. Third, it proposes and evaluates a framework that applies these learned inference techniques to MIS and SC, thereby producing a case study in the effectiveness of this approach. Finally, it situates these contributions within the broader landscape of optimization and learning, highlighting implications for generalization and scalability.

The structure of the thesis reflects this progression from motivation and background to methodology and contribution. Chapter~2 presents the necessary background material, including definitions and examples of combinatorial optimization problems, an overview of machine learning approaches to optimization and a review of Markov Random Fields and their inference algorithms. Chapter~3 discusses the use of MRFs as a general framework for formulating combinatorial optimization problems, illustrating how problems such as Max-Cut, satisfiability, and covering can be expressed as inference tasks. Chapter~4 surveys the integration of machine learning into inference, with emphasis on neural message passing and reinforcement learning methods for combinatorial problems. Chapter~5 introduces the main research contribution of this thesis: the design, implementation, and evaluation of a learned message passing framework for solving MIS and SC. Chapter~6 outlines directions for future research, such as scalability improvements, theoretical guarantees, and potential extensions to other graphical models or quantum-inspired methods. Chapter~7 concludes with a summary of the findings and a reflection on the broader implications of combining machine learning with MRF-based optimization.

\bigskip
